% COLLECTION_METADATA: {"schema_version": 1, "kind": "reused_writeup", "independently_reviewed": false, "primary_source": "attacks/erdos/gpt_pro_5.4/365.tex", "source_paths": ["attacks/erdos/gpt_pro_5.2/365.tex", "attacks/erdos/gpt_pro_5.4/365.tex"], "source_model": "gpt_pro_5.4", "source_completion": 60.0, "source_claim": "unresolved"}
\section*{Erd\H{o}s Problem 365 --- collection record}

\subsection*{PROVENANCE AND STATUS}
Official problem: https://www.erdosproblems.com/365

Saved collaborative-database status: open.
Snapshot checked: 2026-09-23T17:22:24Z.
Latest status and formalization links: https://teorth.github.io/erdosproblems/

This is an attributed collection of existing writeups. The model-folder names below identify their original attribution; placement in GPT\_6\_Astra\_Ultra does not attribute their mathematical work to that model. All locally available source versions selected for this problem are reproduced below, without editing their text.

Proof claims, computations, literature statements, and completion estimates in the reused text have not been independently verified in this collection pass. A source claiming a solution does not change the saved upstream status.

Primary source (latest numbered version in the preferred model folder): \texttt{attacks/erdos/gpt\_pro\_5.4/365.tex}.

\subsection*{REUSED SOURCE 1}
Original attribution: \texttt{gpt\_pro\_5.2}.
Original file: \texttt{attacks/erdos/gpt\_pro\_5.2/365.tex}.

% BEGIN REUSED SOURCE: attacks/erdos/gpt_pro_5.2/365.tex
\section*{Problem 365: Consecutive powerful numbers}

\subsection*{FORMAL RESTATEMENT}
A positive integer $m$ is \emph{powerful} (a.k.a. \emph{squarefull} or \emph{$2$-full}) if for every prime $p\mid m$ we have $p^2\mid m$.

The prompt asks (paraphrasing) whether every pair of consecutive powerful integers $(n,n+1)$ must arise from a Pell-type construction, and then says ``in other words'' whether one of $n,n+1$ must be a square.

\textbf{Ambiguity / minimal correction.} ``Comes from a Pell equation'' is not a mathematically precise property unless one specifies what counts as ``coming from.'' The only precise embedded claim is:

\medskip
\noindent\textbf{(\*)} \emph{If $n$ and $n+1$ are powerful, then at least one of $n$ or $n+1$ is a perfect square.}

I will disprove (\*) by an explicit counterexample.

\subsection*{QUICK LITERATURE/CONTEXT CHECK}
The prompt itself records that many examples come from Pell equations such as $x^2=8y^2+1$ (giving $8y^2$ powerful and $x^2$ a square), but also that Golomb found a counterexample with $12167,12168$, and that Walker produced infinitely many counterexamples via a Pell-type equation $7^3x^2=3^3y^2+1$. (These items are also summarized on the Erd\H{o}s-problems site and in OEIS A060355.)

\subsection*{ATTACK PLAN}
To disprove (\*), it suffices to find one $n$ such that
(i) $n$ is powerful,
(ii) $n+1$ is powerful,
(iii) neither $n$ nor $n+1$ is a square.

\subsection*{WORK (complete counterexample/disproof)}
Take
\[
 n := 12167.
\]
We claim that $n$ and $n+1$ are both powerful but neither is a square.

\medskip
\noindent\textbf{Step 1: $n$ is powerful.}
We have
\[
12167 = 23^3.
\]
Thus the only prime dividing $12167$ is $23$, and its exponent is $3\ge 2$. Hence $12167$ is powerful.

\medskip
\noindent\textbf{Step 2: $n+1$ is powerful.}
Compute
\[
12168 = 8\cdot 9\cdot 169 = 2^3\cdot 3^2\cdot 13^2.
\]
Every prime dividing $12168$ has exponent at least $2$ (namely $3,2,2$), so $12168$ is powerful.

\medskip
\noindent\textbf{Step 3: neither $n$ nor $n+1$ is a square.}
Since $12167=23^3$ has an odd exponent in its prime factorization, it cannot be a perfect square.
Similarly, $12168$ has $2$-adic valuation $v_2(12168)=3$ (odd), so it also cannot be a square.

\medskip
\noindent\textbf{Conclusion.}
The pair $(12167,12168)$ consists of consecutive powerful numbers, but neither term is a square. This contradicts (\*), hence (\*) is false.

\subsection*{VERIFICATION}
We explicitly exhibited complete factorizations:
\[
12167 = 23^3,\qquad 12168 = 2^3\cdot 3^2\cdot 13^2.
\]
Both are powerful because all prime exponents are $\ge 2$. Neither is a square because at least one exponent is odd.

\subsection*{META-CHECK (common failure modes)}
\begin{itemize}
\item \emph{Definition check:} Powerful means every prime exponent is $\ge 2$; both factorizations satisfy this.
\item \emph{Square check:} A square has all exponents even; both numbers have an odd exponent.
\item \emph{Off-by-one:} Verified that $12168=12167+1$.
\end{itemize}

\subsection*{FINAL}
FINAL: LABEL: \textbf{FULL SOLUTION} \quad SUBLABEL: \textbf{COUNTEREXAMPLE/DISPROOF}.

\subsection*{COMPLETION ESTIMATE}
COMPLETION ESTIMATE: 100\%.


% END REUSED SOURCE: attacks/erdos/gpt_pro_5.2/365.tex

\subsection*{REUSED SOURCE 2}
Original attribution: \texttt{gpt\_pro\_5.4}.
Original file: \texttt{attacks/erdos/gpt\_pro\_5.4/365.tex}.

% BEGIN REUSED SOURCE: attacks/erdos/gpt_pro_5.4/365.tex
\section*{Erd\H{o}s Problem \#365 --- ROUND 2}

\subsection*{1) ROUND-2 OBJECTIVE}
\textbf{Path (C): obstruction/correction.}

Round-1 correctly destroyed the first subquestion by one explicit counterexample, but it overclaimed a \emph{full} resolution. The original problem has a second, still-open counting question:
\[
\#\{n\le x : n\mbox{ and }n+1\mbox{ are powerful}\} \stackrel{?}{\ll} (\log x)^{O(1)}.
\]
So the most useful round-2 task is to correct the scope, and at the same time strengthen the first-part disproof from \emph{one} counterexample to an \emph{infinite family}.

\subsection*{2) Round-1 FOUNDATION USED}
I use the following round-1 facts without reproving them.
\begin{itemize}
\item Round-1 isolated the precise first subquestion: if $n$ and $n+1$ are powerful, must one of them be a square?
\item Round-1 gave the explicit counterexample
\[
12167=23^3,\qquad 12168=2^3\cdot 3^2\cdot 13^2,
\]
so the answer to that first subquestion is already \emph{no}.
\end{itemize}

\subsection*{3) NEW INSIGHT / TOOL (ROUND-2)}
The new ingredient is a self-contained Walker-type infinite family of counterexamples coming from the biquadratic unit
\[
\beta:=2\sqrt{7}+3\sqrt{3}.
\]
Odd powers of $\beta^7$ produce infinitely many integer solutions to
\[
7^3x^2=3^3y^2+1,
\]
which in turn give infinitely many consecutive powerful pairs in which \emph{neither} term is a square.

\subsection*{4) ATTACK PLAN (ROUND-2)}
The exact round-1 gap is a \emph{scope gap}, not a proof gap: the overall problem was marked ``full solution'' even though the counting question was untouched.

I will therefore do two things:
\begin{itemize}
\item prove an infinite family of counterexamples to the first subquestion;
\item explain why the overall problem remains unresolved because the counting question survives unchanged.
\end{itemize}

\subsection*{5) WORK (ROUND-2)}
Let
\[
\beta:=2\sqrt{7}+3\sqrt{3},\qquad \beta^{\ast}:=2\sqrt{7}-3\sqrt{3}.
\]
Then
\[
\beta\beta^{\ast}=28-27=1.
\]
A direct expansion gives
\[
\beta^7=2637362\sqrt{7}+4028637\sqrt{3}.
\]
Set
\[
\gamma:=\beta^7=a_1\sqrt{7}+b_1\sqrt{3},
\]
so
\[
a_1=2637362,\qquad b_1=4028637,
\]
and in particular
\[
7\mid a_1,\qquad 3\mid b_1.
\]

For every odd integer $r\ge 1$, write
\[
\gamma^r=a_r\sqrt{7}+b_r\sqrt{3}
\]
with integers $a_r,b_r>0$. (Odd powers have this form because multiplying an odd-form element $u\sqrt7+v\sqrt3$ by an even-form element $A+B\sqrt{21}$ again gives an odd-form element.)

I claim that for every odd $r$,
\[
7\mid a_r,\qquad 3\mid b_r.
\]
This is proved by induction on $r$ in steps of size $2$.
The base case $r=1$ is already known.
Now suppose
\[
\gamma^r=A\sqrt7+B\sqrt3
\]
with $7\mid A$ and $3\mid B$.
Write
\[
\gamma^2=u+v\sqrt{21}
\]
with integers $u,v$. Since
\[
\gamma^2=(a_1\sqrt7+b_1\sqrt3)^2=(7a_1^2+3b_1^2)+2a_1b_1\sqrt{21},
\]
we have
\[
v=2a_1b_1,
\]
so $21\mid v$ because $7\mid a_1$ and $3\mid b_1$.
Now
\[
\gamma^{r+2}=\gamma^r\gamma^2=(A\sqrt7+B\sqrt3)(u+v\sqrt{21})
\]
so
\[
\gamma^{r+2}=(Au+3Bv)\sqrt7+(Bu+7Av)\sqrt3.
\]
Because $7\mid A$ and $7\mid v$, the coefficient $Au+3Bv$ is divisible by $7$.
Because $3\mid B$ and $3\mid v$, the coefficient $Bu+7Av$ is divisible by $3$.
Thus the divisibility property propagates from $r$ to $r+2$, proving the claim.

Next, conjugating in the $\sqrt3$-coordinate gives
\[
\gamma^{\ast}:=a_1\sqrt7-b_1\sqrt3,
\]
and therefore for each odd $r$,
\[
\gamma^r\,\gamma^{\ast r}=1.
\]
Since
\[
\gamma^r=a_r\sqrt7+b_r\sqrt3,
\qquad
\gamma^{\ast r}=a_r\sqrt7-b_r\sqrt3,
\]
we obtain
\[
7a_r^2-3b_r^2=1.
\]
Because $7\mid a_r$ and $3\mid b_r$, the integers
\[
x_r:=\frac{a_r}{7},\qquad y_r:=\frac{b_r}{3}
\]
are well-defined, and the identity becomes
\[
343x_r^2-27y_r^2=1.
\]
Hence the numbers
\[
N_r:=27y_r^2,
\qquad
N_r+1=343x_r^2
\]
are consecutive.
Both are powerful:
\begin{itemize}
\item $27y_r^2=3^3y_r^2$ is powerful because all exponents coming from $y_r^2$ are even and the exponent of $3$ is at least $3$;
\item $343x_r^2=7^3x_r^2$ is powerful for the same reason.
\end{itemize}
Neither is a square:
\begin{itemize}
\item the $3$-adic exponent of $27y_r^2$ is $3+2v_3(y_r)$, which is odd;
\item the $7$-adic exponent of $343x_r^2$ is $3+2v_7(x_r)$, which is odd.
\end{itemize}
So every odd $r\ge 1$ gives a pair of consecutive powerful non-squares.
Since there are infinitely many odd $r$, there are infinitely many such counterexamples.

The first member of this family comes from $r=1$:
\[
x_1=376766,\qquad y_1=1342879,
\]
so
\[
27y_1^2=48689748233307,
\qquad
343x_1^2=48689748233308.
\]
Their factorizations are
\[
48689748233307=3^3\cdot 139^2\cdot 9661^2,
\]
\[
48689748233308=2^2\cdot 7^3\cdot 13^2\cdot 43^2\cdot 337^2.
\]
This is an explicit consecutive powerful pair in which neither term is a square.

Therefore the first subquestion is not merely false; it is false in an infinite family.
However, the second counting question remains. So the original Erd\H{o}s problem is still unresolved overall.

\subsection*{6) ADVERSARIAL VERIFICATION}
\begin{itemize}
\item \emph{Hidden parity issue:} only odd powers of $\gamma$ have the form $a\sqrt7+b\sqrt3$. I used only odd $r$ throughout.
\item \emph{Integrality of $x_r,y_r$:} proved from the inductive divisibility $7\mid a_r$ and $3\mid b_r$.
\item \emph{Could one of the two numbers still be a square if $x_r$ or $y_r$ contains extra factors of $7$ or $3$?} No. The exponents of $7$ and $3$ are respectively $3+2v_7(x_r)$ and $3+2v_3(y_r)$, always odd.
\item \emph{Scope check:} the original problem asks a second counting question, so even an infinite family of counterexamples to the first subquestion does \emph{not} finish the whole problem.
\end{itemize}

\subsection*{7) FINAL}
\textbf{UNRESOLVED (BUT STRICTLY ADVANCED)}

The first subquestion is completely settled negatively, and round-2 strengthens round-1 from a single counterexample to an explicit infinite family. But the overall Erd\H{o}s problem remains open because the counting question is still unresolved.

\subsection*{8) COMPLETION ESTIMATE}
COMPLETION: 60\%

\subsection*{9) REFERENCES}
\begin{itemize}
\item[Bloom365] T. F. Bloom, \emph{Erd\H{o}s Problem \#365}, \texttt{https://www.erdosproblems.com/365}, accessed 2026-03-07.
\end{itemize}

\bigskip\hrule\bigskip


% END REUSED SOURCE: attacks/erdos/gpt_pro_5.4/365.tex

\subsection*{COMPLETION ESTIMATE}
Inherited primary-source estimate: $60\%$. This is the original source's unverified estimate, not a new assessment or confirmation that the problem is solved.
